\chapter{Die Implementierung der IWT}

\section{Einleitung}

Die IWT ist keine rein abstrakte Theorie. Sie wurde numerisch implementiert.

Die Simulation dient drei Zwecken:

\begin{enumerate}
\item \textbf{Validierung:} Überprüfung der internen Konsistenz der Theorie.
\item \textbf{Emergenznachweis:} Demonstration, wie kontinuierliche Physik aus diskreter Dynamik entsteht.
\item \textbf{Bestätigung:} Experimenteller Nachweis, dass die Theorie funktioniert.
\end{enumerate}

Dieses Kapitel beschreibt die Implementierung der Simulation.

\section{Die diskrete Evolutionsgleichung}

Die fundamentale Evolutionsgleichung der IWT lautet:

\[
I_k^{(n+1)} = I_k^{(n)} + T \cdot \Phi_k
\]

mit:

\[
\Phi_k = \sum_{l \in \mathcal{N}(k)} w_{kl}(I_l^{(n)} - I_k^{(n)}) + \lambda \frac{\Delta^2 I_k^{(n)}}{I_k^{(n)}} + \mu I_k^{(n)} \ln \left( \frac{|I_k^{(n)}|}{I_0} \right)
\]

Die drei Terme sind:

\begin{enumerate}
\item \textbf{Lokaler Fluss (Weber):} Diffusion und Glättung der Information.
\item \textbf{Globales Potential (Bohm):} Nicht-lokale Organisation.
\item \textbf{Nichtlineare Strukturbildung:} Lokale Verstärkung.
\end{enumerate}

\section{Die intrinsische Unschärfe}

Aus der diskreten Zeit \( T > 0 \) folgt eine intrinsische Fluktuation:

\[
I_k^{(n+1)} = I_k^{(n)} + T \cdot \Phi_k + \sqrt{\frac{\hbar}{2T}} \cdot \xi_k^{(n)}
\]

Dabei ist \( \xi_k^{(n)} \) eine komplexe, standard-normalverteilte Zufallsvariable.

Dieser Term ist keine externe Störung. Er ist die mathematische Manifestation der bandbegrenzten Frequenzspektrums, das aus der diskreten Zeit folgt.

\section{Die numerischen Parameter}

Die Simulation verwendet folgende Parameter:

\begin{verbatim}
SCALE   = 0.70710678
ALPHA_0 = 1e-2
RHO_0   = 1e-6
RHO_MIN = 1e-8
\end{verbatim}

Die Parameter bedeuten:

\begin{itemize}
\item \texttt{SCALE}: Skalierungsfaktor der Fluktuationen \( \sqrt{\hbar/(2T)} \).
\item \texttt{ALPHA\_0}: Kopplungsstärke des lokalen Flusses.
\item \texttt{RHO\_0}: Referenzdichte für die nichtlineare Strukturbildung.
\item \texttt{RHO\_MIN}: Minimale Dichte (verhindert Singularitäten).
\end{itemize}

\section{Die eingefrorenen Funktionen}

Die Simulation hat zwei eingefrorene Funktionen:

\subsection{Fluktuationen (Vakuum)}

Die Fluktuationen wirken nur im Vakuum.

Sie sind die Quelle der Strukturbildung.

\subsection{Energiesenke (Strukturen)}

Die Energiesenke wirkt nur in Strukturen.

Sie entfernt Energie aus den Strukturen.

\subsection{Die Balance}

Die Balance zwischen Fluktuation und Energiesenke ist gefunden:

\[
\text{Fluktuation} = \text{Energiesenke}
\]

Das System ist stabil.

\section{Die technische Implementierung}

Die Simulation verwendet:

\begin{itemize}
\item \textbf{Knotenzahl:} \( N = 4096 \)
\item \textbf{Iterationen:} 500
\item \textbf{Beschleunigung:} GPU (OpenCL)
\item \textbf{Integrator:} Leapfrog (symplektisch)
\end{itemize}

\section{Die Stabilitätsbedingungen}

Die Simulation ist stabil, wenn:

\[
T < T_{\text{max}} = \frac{r_{\text{min}}}{c}
\]

Die Informationserhaltung wird kontinuierlich überprüft:

\[
\left| \sum_{k=1}^{N} |I_k^{(n+1)}|^2 - \sum_{k=1}^{N} |I_k^{(n)}|^2 \right| < \epsilon
\]

Mit \( \epsilon = 10^{-12} \).

\section{Zusammenfassung}

\begin{itemize}
\item Die IWT wurde numerisch implementiert.
\item Die Evolutionsgleichung enthält drei Terme: lokal, global, nichtlinear.
\item Die intrinsische Unschärfe folgt aus der diskreten Zeit.
\item Die Simulation ist stabil.
\item Die Informationserhaltung wird kontinuierlich überprüft.
\end{itemize}